\section{UX/UI: Design de Interface e Experiência do Usuário}

O design UX/UI da plataforma foi desenvolvido com o objetivo de proporcionar uma interface intuitiva, organizada e de fácil utilização para os produtores de cerveja artesanal durante o acompanhamento do processo de mosturação e realização do teste do iodo.

As telas foram desenvolvidas utilizando prototipação no Figma, permitindo simular o fluxo completo da aplicação antes da implementação do sistema.

A construção da interface priorizou simplicidade operacional, contraste visual e facilidade de navegação, permitindo que o usuário consiga acessar rapidamente as principais funcionalidades da plataforma durante o processo produtivo.

\subsection{Guia de Estilos}

A identidade visual da plataforma foi desenvolvida utilizando cores inspiradas no processo cervejeiro e em interfaces modernas de aplicações web, buscando transmitir organização, contraste visual e facilidade de navegação para o usuário.

A composição visual segue a proporção 60-30-10, técnica amplamente utilizada no design de interfaces para criar equilíbrio visual entre elementos estruturais, componentes secundários e elementos de destaque.

\begin{table}[H]
\centering
\caption{Identidade Visual da Plataforma}
\label{tab:cores}
\begin{tabular}{|l|c|l|p{6cm}|}
\hline
\textbf{Cor} & \textbf{Proporção} & \textbf{Código HEX} & \textbf{Aplicação} \\
\hline

\colorbox[HTML]{FAFAFA}{\rule{0pt}{0.4cm}\hspace{1cm}} &
60\% &
\#FAFAFA &
Cor estrutural da identidade visual, utilizada nos fundos principais da interface e áreas de respiro visual. \\
\hline

\colorbox[HTML]{4E342E}{\rule{0pt}{0.4cm}\hspace{1cm}} &
30\% &
\#4E342E &
Cor secundária aplicada no cabeçalho, menus e fundos secundários da plataforma. \\
\hline

\colorbox[HTML]{F57C00}{\rule{0pt}{0.4cm}\hspace{1cm}} &
10\% &
\#F57C00 &
Cor de destaque utilizada em botões principais, links e chamadas de ação da interface. \\
\hline

\colorbox[HTML]{000000}{\rule{0pt}{0.4cm}\hspace{1cm}} &
-- &
\#000000 &
Cor utilizada para textos principais e elementos informativos da interface. \\
\hline

\end{tabular}
\end{table}

\subsection{Tipografia}

A tipografia da plataforma foi definida com o objetivo de garantir boa legibilidade, organização visual e identidade moderna para a interface do sistema.

As fontes selecionadas foram escolhidas considerando acessibilidade, clareza na leitura e adaptação para diferentes tamanhos de tela.

\begin{itemize}

\item \textbf{Poppins}

Fonte moderna e versátil, utilizada em títulos, chamadas e elementos de destaque da interface. Seus traços suaves e equilibrados transmitem inovação, dinamismo e proximidade visual, contribuindo para uma identidade moderna da plataforma.

\item \textbf{Roboto}

Fonte neutra e altamente legível, utilizada em textos longos e no corpo do conteúdo da aplicação. Sua estrutura proporciona clareza visual e facilidade de leitura em diferentes resoluções e tamanhos de tela.

\end{itemize}